% =========================================================================
% CHAPTER 4: BÁO CÁO TIẾN ĐỘ NGHIÊN CỨU
% =========================================================================

\section{Báo cáo tiến độ nghiên cứu}

% Khung ảnh dùng cho các slide minh họa tiến độ. Khi có ảnh thật, thay lệnh
% \slideScreenshotPlaceholder bằng \includegraphics[...,keepaspectratio]{images/...}.
\newcommand{\slideScreenshotPlaceholder}[2]{%
    \begin{tcolorbox}[
        height=0.65\textheight,
        colback=lightGrayBg,
        colframe=accentTeal,
        boxrule=1pt,
        arc=0pt,
        valign=center
    ]
        \centering
        {\color{accentTeal}\fontsize{28}{30}\selectfont\faImage}\par
        \vspace{0.2cm}
        \textbf{#1}\par
        \vspace{0.12cm}
        {\scriptsize\color{black!60}\texttt{#2}}
    \end{tcolorbox}%
}

% --- SLIDE 4.1: Tổng quan tiến độ ---
\begin{frame}{4.1 Tổng quan tiến độ}
    \begin{table}
        \centering
        \scriptsize
        \setlength{\tabcolsep}{4pt}
        \renewcommand{\arraystretch}{1.28}
        \begin{tabular}{|p{0.33\textwidth}|p{0.18\textwidth}|p{0.42\textwidth}|}
            \hline
            \rowcolor{primaryDark!15}
            \textbf{Khối công việc} & \textbf{Trạng thái} & \textbf{Kết quả hiện tại} \\
            \hline
            Nghiên cứu Campus Network, SDN, SD-WAN & Hoàn thành & Đã hệ thống hóa khái niệm, thành phần và vai trò \\
            \hline
            Phân tích yêu cầu và phạm vi & Hoàn thành & Xác định mô hình đại học gồm Campus chính và ba cơ sở phụ \\
            \hline
            Thiết kế kiến trúc và phân vùng logic & Hoàn thành & Có mô hình phân tầng, VLAN/IP, Server Farm và DMZ \\
            \hline
            Xây dựng mô hình EVE-NG & Đang rà soát & Topology đã hình thành; tiếp tục đối chiếu kết nối và thiết bị \\
            \hline
            Tích hợp các khối SDN và SD-WAN & Đang thực hiện & Hoàn thiện quan hệ điều khiển và kết nối giữa các site \\
            \hline
            Kiểm thử và tổng hợp kết quả & Chưa thực hiện & Thực hiện sau khi mô hình được hoàn thiện \\
            \hline
        \end{tabular}
    \end{table}
\end{frame}

% --- SLIDE 4.2: Mẫu ảnh các Site Campus ---
\begin{frame}{4.2 Minh họa các Site Campus}
    \small
    \begin{columns}[T,onlytextwidth]
        \column{0.52\textwidth}
        \slideScreenshotPlaceholder{CHÈN ẢNH CHỤP SITE}{images/site-campus.png}

        \column{0.44\textwidth}
        \textbf{Cách thiết kế vùng Site:}
        \begin{itemize}
            \item Site 100 dùng mô hình \textbf{Core--Distribution--Access}; các site 200/300/400 được phân vùng theo khoa và VLAN nghiệp vụ.
            \item Mạng người dùng và mạng quản trị được tách thành các subnet riêng theo mã site.
            \item Firewall/WAN Edge tạo ranh giới giữa LAN Campus và mạng WAN.
            \item Prefix của từng site được trao đổi qua SD-WAN overlay.
        \end{itemize}
        \vspace{0.08cm}
        {\color{accentTeal}\textbf{Ảnh cần thể hiện:}} ranh giới site, switch, Firewall, vEdge và uplink về Service Provider.
    \end{columns}
\end{frame}

% --- SLIDE 4.3: Mẫu ảnh Service Provider ---
\begin{frame}{4.3 Service Provider và mạng Underlay}
    \small
    \begin{columns}[T,onlytextwidth]
        \column{0.52\textwidth}
        \slideScreenshotPlaceholder{CHÈN ẢNH INTERNET/MPLS}{images/service-provider-underlay.png}

        \column{0.44\textwidth}
        \textbf{Cách thiết kế vùng Underlay:}
        \begin{itemize}
            \item \textbf{Internet} dùng dải \texttt{203.0.113.0/24}; mỗi WAN Edge nối qua một liên kết transit \texttt{/30}.
            \item \textbf{MPLS} dùng các mạng \texttt{100.64.<site>.x/30} để tách kết nối theo site.
            \item Service Provider chỉ cung cấp IP reachability cho các WAN Edge.
            \item OMP và IPsec hình thành lớp overlay bên trên hai transport.
        \end{itemize}
        \vspace{0.08cm}
        {\color{accentTeal}\textbf{Ảnh cần thể hiện:}} hai cloud Internet/MPLS, các link tới vEdge và Site 900.
    \end{columns}
\end{frame}

% --- SLIDE 4.4: Mẫu ảnh SDN Controller ---
\begin{frame}{4.4 Vùng SDN Controller}
    \small
    \begin{columns}[T,onlytextwidth]
        \column{0.52\textwidth}
        \slideScreenshotPlaceholder{CHÈN ẢNH SDN CONTROLLER}{images/sdn-controller-zone.png}

        \column{0.44\textwidth}
        \textbf{Cách thiết kế vùng điều khiển:}
        \begin{itemize}
            \item Ryu Controller sử dụng \textbf{OpenFlow 1.3} để điều khiển lớp chuyển tiếp.
            \item Control Plane dùng VLAN 99 MANAGEMENT (dải \texttt{10.1.99.0/24}, controller \texttt{10.1.99.10}).
            \item Controller nối 1 cổng vào Server Farm, kênh điều khiển tới 6 switch OVS đi qua link uplink sẵn có.
            \item Lưu lượng điều khiển được tách khỏi Data Plane của người dùng.
        \end{itemize}
        \vspace{0.08cm}
        {\color{accentTeal}\textbf{Ảnh cần thể hiện:}} Controller, 6 switch OVS và kênh control plane trên VLAN quản trị.
    \end{columns}
\end{frame}

% --- SLIDE 4.5: Mẫu ảnh Server Farm ---
\begin{frame}{4.5 Vùng Server Farm}
    \small
    \begin{columns}[T,onlytextwidth]
        \column{0.52\textwidth}
        \slideScreenshotPlaceholder{CHÈN ẢNH SERVER FARM}{images/server-farm-zone.png}

        \column{0.44\textwidth}
        \textbf{Cách thiết kế vùng dịch vụ nội bộ:}
        \begin{itemize}
            \item Server Farm đặt trong \textbf{VLAN 90}, mạng \texttt{10.1.90.0/24}.
            \item DHCP Server \texttt{10.1.90.10} và Syslog Server \texttt{10.1.90.11} dùng địa chỉ tĩnh.
            \item SwitchServerFarm nối dự phòng tới Core-SW1/2 bằng trunk VLAN 90/99.
            \item Vùng dịch vụ được tách khỏi mạng người dùng và DMZ.
        \end{itemize}
        \vspace{0.08cm}
        {\color{accentTeal}\textbf{Ảnh cần thể hiện:}} switch tập trung, DHCP, Syslog và hai uplink Core.
    \end{columns}
\end{frame}

% --- SLIDE 4.6: Mẫu ảnh DMZ ---
\begin{frame}{4.6 Vùng DMZ}
    \small
    \begin{columns}[T,onlytextwidth]
        \column{0.52\textwidth}
        \slideScreenshotPlaceholder{CHÈN ẢNH DMZ}{images/dmz-zone.png}

        \column{0.44\textwidth}
        \textbf{Cách thiết kế vùng dịch vụ công bố:}
        \begin{itemize}
            \item DMZ dùng mạng \texttt{10.1.1.0/28}, tách khỏi Inside và Server Farm.
            \item Web Server \texttt{10.1.1.10} và Mail Server \texttt{10.1.1.11} nối qua SwitchDMZ.
            \item Cặp Firewall ASAv Active/Standby cung cấp gateway và kiểm soát luồng truy cập.
            \item Chỉ các dịch vụ cần công bố mới được phép đi qua chính sách Firewall.
        \end{itemize}
        \vspace{0.08cm}
        {\color{accentTeal}\textbf{Ảnh cần thể hiện:}} Firewall HA, SwitchDMZ, Web/Mail và ranh giới Inside--Outside.
    \end{columns}
\end{frame}

% --- SLIDE 4.7: Kết quả nghiên cứu đã hoàn thành ---
\begin{frame}{4.7 Nội dung đã hoàn thành}
    \begin{columns}[T]
        \column{0.48\textwidth}
        \textbf{Về lý thuyết:}
        \begin{itemize}
            \item Làm rõ mô hình Campus ba lớp và chức năng của từng lớp.
            \item Phân tích nguyên lý tách Control Plane/Data Plane trong SDN.
            \item Phân tích underlay, overlay và các vai trò điều khiển trong SD-WAN.
            \item Xác định vị trí tích hợp SDN và SD-WAN trong kiến trúc chung.
        \end{itemize}

        \column{0.48\textwidth}
        \textbf{Về thiết kế:}
        \begin{itemize}
            \item Xác định bốn site và miền điều khiển SD-WAN.
            \item Hoàn thiện nguyên tắc phân vùng VLAN và quy hoạch địa chỉ.
            \item Phân tách vùng người dùng, quản trị, Server Farm và DMZ.
            \item Xây dựng topology làm cơ sở cho giai đoạn mô phỏng.
        \end{itemize}
    \end{columns}
\end{frame}

% --- SLIDE 4.8: Công việc đang thực hiện ---
\begin{frame}{4.8 Công việc đang thực hiện}
    \begin{block}{Trọng tâm hiện tại}
        Hoàn thiện sự nhất quán giữa cơ sở lý thuyết, thiết kế logic và topology mô phỏng.
    \end{block}
    \vspace{0.15cm}
    \begin{enumerate}
        \item Rà soát vai trò và kết nối của các khối Core, Distribution, Access, Firewall và WAN Edge.
        \vspace{0.12cm}
        \item Đối chiếu VLAN, subnet và vùng dịch vụ với yêu cầu của từng site.
        \vspace{0.12cm}
        \item Hoàn thiện mô hình điều khiển tập trung cho các switch trong Campus chính.
        \vspace{0.12cm}
        \item Hoàn thiện quan hệ underlay/overlay giữa Campus chính và các cơ sở phụ.
        \vspace{0.12cm}
        \item Chuẩn bị các kịch bản kiểm tra kết nối cơ bản cho giai đoạn tiếp theo.
    \end{enumerate}
\end{frame}

% --- SLIDE 4.9: Phần việc tiếp theo trong phạm vi đề tài ---
\begin{frame}{4.9 Kế hoạch hoàn thành giai đoạn hiện tại}
    \begin{enumerate}
        \item \textbf{Hoàn thiện mô hình:} rà soát topology, địa chỉ và quan hệ giữa các thành phần.
        \vspace{0.15cm}
        \item \textbf{Kiểm tra chức năng cơ bản:} kết nối trong cùng VLAN, giữa các VLAN và giữa các site.
        \vspace{0.15cm}
        \item \textbf{Ghi nhận kết quả:} thu thập sơ đồ, bảng trạng thái và minh chứng từ mô phỏng.
        \vspace{0.15cm}
        \item \textbf{Hoàn thiện báo cáo:} đồng bộ phần lý thuyết, thiết kế và tiến độ thực hiện.
    \end{enumerate}
    \vspace{0.2cm}
    \begin{block}{Giới hạn của kế hoạch}
        Chỉ tập trung hoàn tất các hạng mục đã xác định trong đề tài, không mở rộng thêm công nghệ hoặc phạm vi mới.
    \end{block}
\end{frame}
